\section{A locally evolving outer representation}\label{sec:outer-representation}

We now combine the macrograph, the outer codes, and the certified programs
into one representation. A long preparation is part of the represented
path: its current position is recorded by a binary clock, while its scratch
space is reused. Thus preparation time does not multiply the length of a
represented vertex. The main additional requirement is that every retained
certificate, including certificates introduced to verify other
certificates, can change locally when a logical tape entry changes.

All notation in this section belongs to the outer construction. In
particular, $J$, $H$, and $\Gamma=\F_2^r\rtimes H$ are those of
Section~\ref{sec:outer-codes}. We retain $M_o=2^r$, the tape length $h$,
the four logical arrays $b$, and the persistent tuple $P(b)$ of
Lemma~\ref{lem:outer-wiring}. The assigned degree $k_S$ of a primary
$F_S$ counts its variable factor occurrences, and $d_2$ is the fixed upper
bound on the number of factor occurrences. In particular,
\[
 r=O(\log N),\qquad M_o=O(N),\qquad |H|\leq r^9,
 \qquad M_o>(8d_2+8)h.
\]
Every persistent program stage has a specified linear payload recipe in
the primaries. This includes support masks, transforms, routing checks,
remainder computations for inactive phases, and every matrix-valued
reencoding intermediate. Its assigned degree is at most the maximum
assigned degree of the primary inputs to that recipe.

\subsection{Finite differences of the complete certificate tuple}

Fix a start frame $c\in\{0,\ldots,4h-1\}$. The finite edit alphabet
$\Sigma$ of Lemma~\ref{lem:macrograph} describes additive changes
$h_\sigma(c)$ of the logical arrays. Enlarge this fixed alphabet by inverse
edits and a zero edit if necessary. All addresses depend only on
$(c,\sigma)$ and public routing data. The choice of the edit can depend on
the critical logical bits; a recipe for a specified edit cannot.
For a field-valued function $F$ define
\[
 \Delta_\sigma F(b)=F(b+h_\sigma(c))-F(b).
\]
This definition is made on arbitrary field arrays $b$, independently of
whether they satisfy the macrograph conditions.

For a variable singleton occurrence $a$, routing gives either no change or
a change $\epsilon_{a,\sigma}L_{j_{a,\sigma}}(Z)$, where
$\epsilon_{a,\sigma}\in\{-1,0,1\}$ and $j_{a,\sigma}<h$ are public at
this frame. A default address is used when $\epsilon_{a,\sigma}=0$.
Public factors have zero increment. If $S_{\rm var}$ is the set of
variable \emph{occurrences} in $S$, distributivity gives
\begin{equation}\label{oe:primary-difference}
 \Delta_\sigma F_S
 =\sum_{\varnothing\ne A\subseteq S_{\rm var}}
     F_{S\setminus A}(b)
     \prod_{a\in A}\epsilon_{a,\sigma}L_{j_{a,\sigma}}.
\end{equation}
Repeated occurrences are retained separately in this formula. Every
primary on its right has assigned degree at most $k_S-1$. Multiplication
by a specified Lagrange polynomial is the linear coefficient program of
Lemma~\ref{op:scalar-programs}. The number of subsets in
\eqref{oe:primary-difference} is bounded in terms of $d_2$ alone.
Every ordinary polynomial product appearing here has degree at most
$d_2(h-1)<M_o$; consequently its implementation by cyclic coefficient
convolution has no wraparound error. All possible addresses and zero
coefficients use a common padded program topology.

\begin{lemma}[Closure under finite differences]\label{oe:derivative-closure}
For every positive-degree persistent full slot $D$ and every
$\sigma\in\Sigma$, one can allocate a finite certified full program whose
output, in the same encoder type as $D$, is the increment of $D$ under
$b\mapsto b+h_\sigma(c)$. One can do this simultaneously for every
positive-degree full intermediate introduced by these programs, including
all their verification intermediates. The resulting list has
$\operatorname{poly}(r)$ slots and description bits. Every new slot has a
specified linear recipe in persistent primaries of strictly smaller
assigned degree than its parent target. All identities hold on arbitrary
field arrays $b$.
\end{lemma}

\begin{proof}
For a primary use \eqref{oe:primary-difference}. If $D$ is another
persistent stage with linear payload recipe $\mathcal L_D$, apply
$\mathcal L_D$ to the corresponding primary increments, replacing
degree-zero primary increments by zero. At this fixed frame its
coefficients, pulls, and encoder types are fixed. Hence
\[
 \mathcal L_D\bigl((F_S(b+h_\sigma(c)))_S\bigr)
 -\mathcal L_D\bigl((F_S(b))_S\bigr)
 =\mathcal L_D\bigl((\Delta_\sigma F_S(b))_S\bigr).
\]
The operations used in Section~\ref{sec:outer-programs} implement this
identity at the payload level. In particular, multiplying a scalar slot
by the public matrix $\rho_e(k)$ is linear in that scalar; changing one
encoder axis or contracting matrix entries is also linear. These
operations introduce no additional logical degree.

Allocate full slots for every stage of each increment program, with new
slots for its mixed and untwisted verification stages. An identity or zero
stage may be inserted to give an output its required type. Order stages
topologically, so that Lemma~\ref{lem:outer-fill} can initialize one full
target using already completed parents. For every newly allocated stage,
record the ancestor directed acyclic graph that expresses its payload in
the persistent primaries. Its roots have degrees at most one less than
those of the target whose increment is being computed. If its assigned
degree is positive, apply the same construction to this new stage for
every edit. This includes a stage introduced solely for verification: it
will be retained at full preparation and therefore needs its own
increment.

Inlining here means substitution in a directed acyclic graph, with
sharing retained within each target computation. If a primary is used
many times, substitute its increment program once and fan out its result.
Separate target computations may use fresh copies. Along any chain of
newly generated targets the assigned degree drops by at least one.
Therefore the construction stops after at most $d_2$ derivative tiers.
At each tier the number of target graphs, their sizes, their ancestor
descriptions, and their substitutions are bounded by a fixed polynomial
in the preceding total description size and $r$. A fixed number of such
polynomial expansions remains polynomial in $r$.

The programs verify the computed increments by their own stage
relations, starting from the lower-degree primaries. They do not test a
finite-difference assertion about an otherwise unspecified witness.
Neither their construction nor their identities use the final validity
equations. This also ensures that increment recipes remain defined on
field data produced by an edit that is not a legal macrograph move.
\end{proof}

\subsection{Public addresses, logical reads, and a reversible preparation}

Public address tables need not admit succinct evaluation at a hidden
frame $c$. We therefore represent the needed access explicitly. For each
field entry $a(c)$ of the public data table, form the scalar payload
\begin{equation}\label{oe:public-walsh-table}
 V_a(x,u)=\sum_{0\leq c'<4h}a(c')(-1)^{x\cdot c'}.
\end{equation}
Integers in this display are identified with their $r$-bit vectors, and
the coefficient table is zero outside $[0,4h)$. The payload is independent
of $u$. Hardwire its full encoding at every publicly requested symbol
position. The table includes operand addresses, rule types,
terminal-comparison addresses, and choice-rule data for both directions,
together with the
permuted edit addresses of every factor type used by
Lemma~\ref{oe:derivative-closure}. Addresses for inactive phases are
included too. Strings and indices are split into field entries with fixed
parsing conventions. There are only $\operatorname{poly}(r)$ entries per
frame.

The first preparation steps are an A-only public focus. For the bits of
the repeated frame $c$, apply the mode selectors
\[
 \frac{I+(-1)^{c_i}R_{(e_i,1)}}2
\]
successively, implementing a nonfirst-bit pull by the fixed generators
as in Section~\ref{sec:outer-programs}. Applied to
\eqref{oe:public-walsh-table}, their completed output is
$a(c)(-1)^{x\cdot c}$ on the systematic payload. Removing that known sign
recovers $a(c)$ identically at every center. No recurrence using a
table-dependent coefficient is active before the public focus is
complete. Retain this focus until preparation has been reversed. On
malformed data use fixed default field values, indices, and instructions;
thus every local algorithm is total.

Next, in the topological order of Lemma~\ref{oe:derivative-closure}, fill
every transient full computation stage by Lemma~\ref{lem:outer-fill}.
Previously completed targets remain full and their temporary fill scratch
is empty. At every intermediate location check their direct full-operation
relations; apply the two-buffer fill checks to the current target; and
check zero in all still empty full targets and inactive scratch slots.
The persistent tuple $P(b)$ is checked throughout. All these conditions
are tested together at the same center.

Finally prepare the critical logical reads. The persistent slot $U_D$
provided by Lemma~\ref{lem:outer-wiring} has payload
\[
 U_D(x)=\sum_{j<h}(-1)^{x\cdot j}D_j.
\]
Its Walsh coefficient at address $j$ is $D_j$. Using the completed public
address focus, build A-only selector stacks for the fixed list of
critical roles needed by either the forward or reverse macro proposal
and its existence flag. Unused roles use default addresses. Include
enough roles for each logical array type; their number is constant.
Build the identical focus stacks on the full increment outputs for
$U_D$ for every $\sigma$. The recipe for $U_D$ uses just one degree-one
singleton, so each increment input has assigned degree zero. At a fixed
frame its entire focused stack is independent of $b$ as well.

\begin{lemma}[Preparation and local evolution]\label{oe:preparation}
There is an integer $J_{\max}\geq1$ and a canonical preparation indexed by
$j=0,\ldots,J_{\max}$ with the following properties.
\begin{enumerate}[label=\textup{(\roman*)}]
\item At $j=0$ all frame-dependent full and A-only transient slots are
zero. At $j=J_{\max}$ all full increment programs and critical focuses are
complete, and every full-target fill buffer and orbit-focus buffer is
zero.
\item For every field array $b$ and frame $c$, every snapshot is uniquely
specified. Its transition to the next or previous snapshot uses a fixed
finite list of permutation-neighbors of each requested center.
\item The number of steps is $J_{\max}=O(M_o\operatorname{poly}(r))$.
The number of stored slots, the schedule description, and the work to
interpret one microclock or perform one local transition are
$\operatorname{poly}(r)$, including entry dimensions and bit lengths.
\item At full preparation, applying an edit $h_\sigma(c)$ changes the
entire retained tuple to its canonical value on $b+h_\sigma(c)$ by local
addition of the stored increment outputs. The public focus is unchanged.
\end{enumerate}
\end{lemma}

\begin{proof}
Concatenate the public focus, all full-target filling schedules, and the
critical-read focuses. The stage lists for selectors and cutoffs have
fixed padded lengths independent of their address values. The
per-target mode schedule consists of $M_o$ iterations of fixed
polylogarithmic sublength, together with its initial and final buffer
steps. The exceptional first iteration and final erasure are specified
separately. Thus a binary microclock can be interpreted by scanning the
polynomial-in-$r$ list of program boundaries and performing arithmetic
on $O(\log M_o+\log\operatorname{poly}(r))$ bits. The interpreter selects
the current iteration, buffer, initialized prefix, and transition; it
does not execute all preceding iterations. Distinct programs may use
distinct scratch indices to keep their shapes and twists fixed.
An idle preparation step supplies $J_{\max}\geq1$ if needed.

For the last assertion, write $D_\sigma(b)$ for the stored full increment
of a positive-degree full slot $D$. At the top perform, simultaneously,
\begin{equation}\label{oe:simultaneous-edit}
 D(b)\longmapsto D(b)+D_\sigma(b)
             =D(b+h_\sigma(c))
\end{equation}
for every persistent and transient full computation slot. Equation
\eqref{oe:simultaneous-edit} is valid in its own encoder because that
encoder is linear. The increment slots themselves are included in this
update whenever their assigned degree is positive; their required
increments exist by Lemma~\ref{oe:derivative-closure}. Degree-zero slots
do not change. All additions read the old tuple before writing the new
one.

The full-target temporary buffers are already empty. Public focus
depends only on $c$ and remains fixed. Each main critical-read stack is a
linear selector-and-pull program at fixed public addresses. Adding the
matching focused increment stack therefore updates all its initialized
stages to their new canonical values. Those increment stacks have degree
zero and stay fixed. This verifies the assertion also for the retained
A-only data. The preparation can now be erased in the reverse order,
using exactly the same recipes on the new $b$ and unchanged frame.
\end{proof}

\subsection{The extended graph and its endpoint payload}

Define an extended valid label to be
\begin{equation}\label{oe:extended-label}
 v=(c,b,\eta,j),\qquad
 0\leq c<4h,\quad \eta\in\{0,1\},\quad 0\leq j\leq J_{\max},
\end{equation}
where $b$ is a valid macrograph state at stage $c$ if $\eta=0$, and at
stage $c+1\pmod{4h}$ if $\eta=1$. We call these modes pre and post,
respectively. Its raw word $R_v$ contains, at every $k\in\Gamma$, all four
parts of the persistent tuple $P(b)$, all full and A-only transient data
in the canonical snapshot $(c,b,j)$, and a canonically encoded repeated
copy of $(c,\eta,j)$ and its range information. Padding conventions will
be fixed below.

In pre mode increase $j$ until the top. At the top use the prepared
critical reads to compute the macrograph's forward choice and existence
flag. If the edge exists, perform \eqref{oe:simultaneous-edit}, leave
$c,j$ fixed, and change the mode to post. In post mode decrease $j$ until
zero, then perform
\begin{equation}\label{oe:frame-advance}
 (c,b,1,0)\longmapsto(c+1\bmod4h,b,0,0).
\end{equation}
The reverse rules use the reverse preparation transition, the reverse
macro choice at the same start frame, and the inverse of
\eqref{oe:frame-advance}. Both directions read their choice and flag
before the top update. At a missing edge the corresponding proposed word
is defined to be the current raw word and its flag is zero. This is a
total proposal convention, not a graph self-edge. At a malformed input
the local proposal and flag likewise have a fixed total default.

For each valid macro vertex at stage $s$, the extended graph contains the
directed path
\begin{equation}\label{oe:replacement-path}
 \begin{split}
 (s-1,b,1,J_{\max})&\longrightarrow\cdots
 \longrightarrow(s-1,b,1,0)\\
 &\longrightarrow(s,b,0,0)
 \longrightarrow\cdots\longrightarrow(s,b,0,J_{\max}),
 \end{split}
\end{equation}
with frames taken modulo $4h$. Macro edges join the corresponding ends
of these paths. Every internal edge changes the clock or mode, and every
top edge changes mode, so none is a self-edge. The graph has indegree and
outdegree at most one. Every label belongs to a path of positive length;
there are no isolated vertices.

Let $(s_*,b_*)$ be the public exceptional macro source of
Lemma~\ref{lem:macrograph}. The unique exceptional extended source and
its actual successor are
\begin{equation}\label{oe:public-source}
 v_*=(s_*-1,b_*,1,J_{\max}),\qquad
 v_*^+=(s_*-1,b_*,1,J_{\max}-1).
\end{equation}
Their entire raw words are deterministically computable in
$\operatorname{poly}(N)$ time. Indeed the macro source is public, all
full encoders and program formulas are computable by global array
algorithms, and even explicitly following this one preparation takes
only polynomially many steps.

An extended endpoint can occur only at a top label of
\eqref{oe:replacement-path}. Map it to its macro vertex by retaining $b$
and taking stage $c+\eta\pmod{4h}$. If it is not $v_*$, the endpoint
decoder of Lemma~\ref{lem:macrograph} returns a nonzero string $y$ with
$P(S(y))\ne y$ or $S(P(y))\ne y$. Thus the endpoint payload consists of
the logical tapes and the repeated label; it does not require finding a
history of the transient computation.

\subsection{A packed tester for valid represented vertices}

Let $\Pi$ contain identity, all forward and inverse A- and B-neighbor
maps needed by the codes, the left $H$ generators used for payload
independence, the right first-bit shifts and single-cycle rotations and
their inverses, and the right $H$ generators used by scratch programs.
Add the finitely many bounded compositions required by a primitive local
expression. Its cardinality $q=|\Pi|$ is fixed independently of $N$.
Arbitrary powers of a rotation or paths through $H$ are represented by
successive stages; they do not add new query maps to this list. At every
center the tester reads $W|_{\Pi(k)}$ and rejects if \emph{any} applicable
condition below fails:
\begin{enumerate}[label=\textup{(\roman*)}]
\item canonical parsing, legal repeated clock, range checks, and zero
padding;
\item all full-slot local codes, A/B agreements, and systematic-payload
duplication, including targets that are partial or inactive;
\item all A-only local codes and A agreements, including inactive
scratch;
\item all structural persistent relations and the macro-validity
conditions selected by $(c,\eta)$;
\item every required current-snapshot relation: completed full-operation
checks, the current partial-target comparison, initialized scratch-prefix
recurrences in both buffers and all focuses, and inactive-zero checks.
\end{enumerate}
In particular, the tester does not choose one relation uniformly from a
growing list. All relations share this fixed set of center permutations.
Write $\operatorname{rej}(W)$ for its rejection probability at uniform
$k\in\Gamma$, and use normalized whole-symbol distance
\[
 \dist(W,W')=\Pr_{k\in\Gamma}[W(k)\ne W'(k)].
\]

\begin{lemma}[Soundness of the valid-word tester]\label{oe:valid-testing}
There are fixed constants $C_{\rm rep}<\infty$ and $g_{\rm rep}>0$
such that, if $\operatorname{rej}(W)<g_{\rm rep}$, some valid extended
label $v$ satisfies
\[
 \dist(W,R_v)\leq C_{\rm rep}\operatorname{rej}(W).
\]
The tester has perfect completeness. Conversely,
$\operatorname{rej}(W)\leq q\dist(W,R_v)$ for every valid label $v$.
Distinct valid labels have raw distance at least a fixed
$\delta_{\rm raw}>0$.
\end{lemma}

\begin{proof}
Put $\theta=\operatorname{rej}(W)$. The parsing, repetition, and full-code
proximity conclusions of Theorem~\ref{thm:outer-codes} first change
$O(\theta)$ centers to give legal constant clock data and exact full
codewords. A corrected repeated clock of illegal range would cause
rejection everywhere, so small enough $\theta$ rules it out. Exact
one-axis distance then forces systematic-payload duplication in every
full slot. Otherwise that one failed duplication relation already has a
fixed positive rejection probability. Zeroing locally A-invalid scratch
parts costs another $O(\theta)$ centers. A change on $z$ centers changes
any simultaneous test on at most $qz$ centers; this estimate uses the
common permutation list and does not depend on the number of slots.

We next restore all scratch relative to these corrected full codewords
and the corrected clock, before using any advertised full-operation
identity. The reference for the public focus is the actual hardwired
public word at the corrected frame. Treat its recurrences first. A later
table-dependent recurrence is active only when this reference focus is
complete and hence gives identical correct coefficients at every center.
Each initial expression $X$ in a mode-fill program is A-coded from its
exact full parents, even if those parents do not yet satisfy their
advertised payload recipes. The critical and increment focuses also
start from exact A-coded full words. Thus every required reference
prefix is exactly A-coded.

These are precisely the reference and ordering hypotheses of
Lemma~\ref{op:scratch-restoration}. The preprocessing above changed only
$O(\theta)$ sites, so the common query list bounds both the joint
A-conflict fraction and the simultaneous recurrence-failure fraction of
this one preprocessed word by $O(\theta)$. Apply the lemma to its entire
scratch tuple, including inactive zeros. Its fixed-word induction keeps
every prefix estimate $O(\theta)$ and makes only one final replacement.
All scratch is therefore restored at cost $O(\theta)$, independently of
the number of stages. No full-operation identity has been used to obtain
this conclusion.

Call the resulting tuple $\widetilde W$. Its remaining tests still fail
with probability at most $C'\theta$ for a fixed $C'$. By
Lemma~\ref{lem:outer-wiring}, either its persistent tuple is exactly
$P(b)$ for logical bits $b$ valid at the selected macro stage, or a
persistent relation fails with fixed positive probability. Choose
$\theta$ below that gap divided by $C'$ to obtain the former alternative.
Now follow the initialization order of the transient full targets.
Completed targets with canonical parents are forced to their exact
outputs by the shared-axis gap following Definition~\ref{op:full-operation}.
For the current partial target, its required cutoff comparison has the
same gap by Lemma~\ref{lem:outer-fill}. The cutoff scratch has already
been restored: its completed value is precisely $P_{t'}X$. In the case
that the target shares only the B axis with $X$, orbit invariance makes
this completed reference B-coded. Only at this point is B-axis distance
applied. No intermediate A-only scratch is assumed to be B-valid.
Inactive full targets are forced to zero by full-word distance.

There are only finitely many kinds of these gaps, uniformly in their
stage indices, program lengths, and vector dimensions. Their minimum is
positive. If any relation failed exactly on $\widetilde W$, it alone
would give at least that gap in the packed test. Decrease $\theta$ once
more to rule out every such failure. All full slots and restored scratch
then have the unique canonical values of $(c,b,\eta,j)$. The total
restoration cost is $C_{\rm rep}\theta$ for a fixed $C_{\rm rep}$.

Canonical snapshots satisfy every test by their construction, proving
perfect completeness. If a neighborhood avoids the disagreement set
between $W$ and $R_v$, its test accepts; a union bound over $\Pi$ proves
the converse. Finally, two labels with different repeated data differ
at every center. If only their logical arrays differ, one logical
singleton interpolant differs. Its full encoded difference is nonzero
and has the fixed distance of Theorem~\ref{thm:outer-codes}, including the
constant conversion between four-part and whole-symbol metrics. This
gives $\delta_{\rm raw}>0$.
\end{proof}

\subsection{Full representation cost and a fixed local Boolean exponent}

We record the size argument before choosing any correction accuracy or
any number of Boolean ports. This order is needed later. All constants
in the following count depend only on the fixed gate basis, $d_2$, and
the already chosen outer code parameters.

One-bit alignment takes $O(r)$ elementary rotations. A conjunction of
bit conditions takes $O(r^2)$ primitive stages, and a carry or borrow
shift takes $O(r^3)$. The Lagrange window product takes $O(r^4)$ stages.
The other transforms and masks in Section~\ref{sec:outer-programs} fit
within this last bound. There are $O(r)$ initial factor and row families:
a fixed number of row lists, with $O(r)$ routing layers for each needed
permutation and a bounded subset list for each monomial. Hence there are
$O(r^5)$ persistent primitive stages. Storing their ancestor graphs,
parent indices, types, and coefficient instructions is bounded, with
ample slack, by $Cr^{20}$ description bits. Public indexed field data
are counted separately as local constant words below; descriptors refer
to their entries by indices, rather than copying an entire length-$h$
array into a descriptor.

Here is an explicit finite-depth bound for derivative expansion. Let $T$
bound the complete current list of graph descriptions and subgraph
descriptions, with each slot, parent occurrence, and instruction counted
at least once. A target graph has at most $T$ nodes and at most $T$
primary roots. Substitution at one root uses a bounded number of the
$O(r^4)$ Lagrange programs. There are at most $T$ target graphs. Recording
an ancestor subgraph for every newly introduced stage, retaining sharing
within it and allowing fresh copies for separate targets, costs at most
the square of the size of its target computation. Public address
instructions and index encodings contribute fixed powers of $r$ and
logarithms of these list lengths. The following deliberately loose bound
therefore covers all these operations at one tier:
\begin{equation}\label{oe:recipe-size-recurrence}
 T_{a+1}\leq C r^{20}(T_a+1)^6,\qquad T_0\leq Cr^{20}.
\end{equation}
For example, the six factors allow a separate factor for choosing a
target, traversing its graph, substituting its roots, and copying the
ancestor descriptions, with two additional factors for indexing and
duplication; the actual substitution keeps more sharing. If
\begin{equation}\label{oe:recipe-exponent}
 E_0=20,\qquad E_{a+1}=6E_a+20,\qquad E=E_{d_2},
\end{equation}
then all descriptions are $O(r^E)$. This recurrence is composed only
$d_2$ times, by Lemma~\ref{oe:derivative-closure}.

For each full target, a cutoff program uses at most $|H|$ group elements,
$r$ prefix rectangles, $r$ selectors per rectangle, and $O(r)$ primitive
alignment stages per selector. This is $O(r^{12})$ stages. For the orbit
focus, there are at most $|H|$ frequencies and $|H|$ requested group
indices. Even finding and storing a separate generator path of length
$|H|$ for every such pull gives $|H|^3$ stages, in addition to the mode
selectors. A bound $O(r^{30})$ covers both focus and cutoff work and the
two buffers. An entry alphabet has dimension at most $|H|^2\leq r^{18}$,
and a field entry uses $O(r)$ bits. Therefore
\begin{equation}\label{oe:raw-size-bound}
 \text{raw data, local public constants, and schedule metadata at one
 center}\ \leq C r^{E+60}.
\end{equation}
This includes public-address focus, critical focus, all slot types and
all four parts. The long sequence of $M_o$ cutoffs is represented by a
clock and reusable buffers, not by $M_o$ copies of these data.

Put $\ell_N=\lceil\log_2(N+2)\rceil$. Fix once an integer $C_L>E+100$
and pad each raw symbol to
\begin{equation}\label{oe:raw-length}
 L=(\ell_N+C_{\rm base})^{C_L}.
\end{equation}
A fixed lower bound on $C_{\rm base}\geq1$ makes
\eqref{oe:raw-size-bound} fit for every $N$. The unpadded construction is
still parameterized by $r$ and the same slot schema. Increasing
$C_{\rm base}$ later adds zero padding only. The raw word and all
center-specific public descriptions have total length
\begin{equation}\label{oe:total-raw-size}
 |\Gamma|\operatorname{poly}(\ell_N)
 =M_o|H|\operatorname{poly}(\ell_N)
 =N\operatorname{poly}(\ell_N).
\end{equation}
This is a bound on the listed bits, including index and field
representations, rather than just the number of slots.

Every local test and proposal can now be implemented by a fixed
polynomial-time bit algorithm on a bounded number of $L$-bit symbols.
Parsing, loops over slots and the explicit group $H$, arithmetic in $J$,
program interpretation, and multiplexing within a supplied
polylogarithmic array use bounded nests of polynomially bounded loops.
The largest local computation is the orbit-restricted parity solve of
Lemma~\ref{lem:outer-fill}. Its dimension is at most
\[
 C|O|\,|H|\,\dim_J(U)\leq C|H|^4\leq CL^4.
\]
Its explicitly stored matrix has $O(L^9)$ bits if each field entry is
bounded by $L$ bits. Gaussian elimination uses $O(L^{12})$ field
operations, and elementary arithmetic and inversion have a schoolbook
bound $O(L^4)$ in this generous bit-length allowance. Even accessing
each entry by a linear scan through the matrix, and scanning the full
slot and metadata lists, adds only fixed polynomial factors. All other
displayed local operations have the same type of finite loop bound.
No routine iterates through the $M_o$ cutoffs to answer one query.
Choose an integer $a$ exceeding the sum of these fixed loop and
arithmetic exponents. A standard direct time-unrolling of a deterministic
bit computation of length $O(L^a)$ gives an acyclic bounded-fanin Boolean
circuit of size $O(L^{2a+2})$: represent each configuration by its used
tape cells and head indicators, and implement each next configuration
by Boolean selection. Thus one fixed exponent $c_{\rm raw}$ bounds all
raw local circuits. Its precise numerical value is immaterial; the
displayed construction fixes it before correction rounds or port arities
are selected.

All lists and constants are uniformly constructible. Global algorithms
enumerate the finite fields, groups, public routing tables, and full
constant encodings, using the deterministic searches and linear algebra
of Sections~\ref{sec:outer-codes} and~\ref{sec:outer-programs}. Their
time is $\operatorname{poly}(N)$, though it need not be quasilinear.
Input-dependent data without succinct evaluation at $c$ are accessed
through \eqref{oe:public-walsh-table}; other public constants are
hardwired only at the bounded, publicly specified neighbors of the
requested center. There is no hidden access to an unrepresented
length-$N$ array.

\subsection{Correction and the complete outer interface}

All full and scratch slots of $R_v$ satisfy their A constraints. Apply
the joint A-axis corrector of Lemma~\ref{lem:outer-uncertainty} to all
of them, together with majority correction of the repeated clock and
canonical padding. The local decoder enumerates at most $2^d$ label
subsets and solves the common scalar constraints componentwise. Since
$d$ is fixed, this again has a fixed polynomial bound in $L$, independent
of the number or dimensions of the packed slots. Enlarge $c_{\rm raw}$
once to cover this one-round routine.

Let $Q\geq2$ bound its neighborhood size and the sizes of the raw
neighborhoods. For any fixed nonnegative integer $t$, unrolling $t$
correction rounds followed by a raw local routine uses at most
\begin{equation}\label{oe:fixed-interface-composition}
 C(1+Q+\cdots+Q^t)=O(Q^t)
\end{equation}
copies of these same $L$-bit circuits. Their intermediate outputs still
have length $L$. It follows that, for a constant $C_t$ depending on $t$,
the composite circuit has size at most $C_tL^{c_{\rm raw}}$. Its ports
are the finitely many compositions of the original permutations to depth
at most $t$, followed by a raw neighborhood. The exponent is independent
of $t$. A later increase of $C_{\rm base}$ changes only the linear work
of checking and writing padding; it changes neither this exponent nor
the number of ports. In particular, fixed multiplicative constants may
be absorbed into a subsequent exponent allowance with slack by increasing
$C_{\rm base}$ after all those constants have been fixed.

\begin{theorem}[Outer representation interface]\label{thm:outer-representation}
There are fixed positive constants $\delta_{\rm raw},\rho_{\rm raw}$,
fixed integers $C_L,c_{\rm raw}$, and uniform deterministic algorithms
with the following properties for every End-of-Line input of length $N$.
They construct a directed graph of valid extended labels, represented by
canonical words
\[
 R_v\in(\{0,1\}^L)^\Gamma,
 \qquad |\Gamma|=N\operatorname{poly}(\ell_N),
 \qquad L=(\ell_N+C_{\rm base})^{C_L},
 \qquad \ell_N=\lceil\log_2(N+2)\rceil,
\]
where $C_{\rm base}$ may be any sufficiently large fixed integer.
\begin{enumerate}[label=\textup{(\roman*)}]
\item \textbf{Graph and endpoint data.} The graph has indegree and
outdegree at most one, no self-edges, and no isolated vertices. It has the
public distinguished source $v_*$ and actual successor $v_*^+$ in
\eqref{oe:public-source}; their entire words are computable in
$\operatorname{poly}(N)$ time. Every other endpoint is mapped in
$\operatorname{poly}(N)$ time to a nonzero solution of the input
End-of-Line instance. Every canonical word is computable from its valid
label in polynomial time.
\item \textbf{Distance and valid-word testing.} Distinct canonical words
have normalized symbol distance at least $\delta_{\rm raw}$. A uniform
center test has perfect completeness and uses only a fixed finite list
$\Pi$ of permutation-neighbors, of cardinality $q$ independent of $N$.
For every fixed $\alpha>0$ there is a fixed $\theta_\alpha>0$ such that
\[
 \operatorname{rej}(W)<\theta_\alpha
 \quad\Longrightarrow\quad
 \dist(W,R_v)<\alpha\quad\hbox{for some valid }v.
\]
Also $\operatorname{rej}(W)\leq q\dist(W,R_v)$ for every valid $v$.
\item \textbf{Local proposals and flags.} At each requested center,
total deterministic routines on $\Pi(k)$ compute a next or previous
proposed symbol and an existence bit. On $R_v$ the bit is independent
of the center and equals the actual edge-existence flag. If it is one,
the proposed symbols form exactly the canonical neighboring word; if it
is zero, the default proposal is $R_v$ itself.
\item \textbf{Value-independent uncertainty correction.} There is a
total local one-round corrector with a bounded permutation neighborhood.
If $|U|/|\Gamma|\leq\rho_{\rm raw}$, it maps every word agreeing with
$R_v$ outside $U$ to a word agreeing with $R_v$ outside a set $U^+$
determined only by $U$, with $|U^+|\leq|U|/2$. Consequently $t$ rounds
leave a prescribed uncertainty support of density at most
$2^{-t}|U|/|\Gamma|$. The guarantee holds for arbitrary values on $U$.
Different requested-center simulations may use different values there,
provided repeated occurrences of a position are consistent within each
individual cone.
\item \textbf{Global recovery.} From any word within
$\rho_{\rm raw}$ symbol distance of a valid word, a deterministic
polynomial-time algorithm recovers its exact extended label and raw
word. It then recovers the logical bits by systematic coefficient
extraction and evaluation on the subgroup. In particular it extracts an
End-of-Line answer whenever the recovered label is a nonexceptional
endpoint.
\item \textbf{Uniform cost and order of choices.} The total raw data and
listed public descriptions have $N\operatorname{poly}(\ell_N)$ bits
and are constructible in $\operatorname{poly}(N)$ time. Each raw local
test, proposed symbol, flag, and one-round corrector has a Boolean
circuit of size at most $CL^{c_{\rm raw}}$ on a fixed number of symbols.
For each fixed $t$, correction for $t$ rounds followed by any such raw
routine has size at most $C_tL^{c_{\rm raw}}$ and at most $O(Q^t)$
symbol ports, with uniform construction for that $t$. The exponents
$C_L,c_{\rm raw}$ are fixed before $t$ and any later port arity. Further
fixed additive padding in $C_{\rm base}$ has only linear processing
cost and does not change these exponents or arities.
\end{enumerate}
\end{theorem}

\begin{proof}
The graph and local evolution were established in
Lemma~\ref{oe:preparation} and \eqref{oe:replacement-path}. The tester
and distance are Lemma~\ref{oe:valid-testing}. For example, its constants
allow
\[
 \theta_\alpha=
 \min\{g_{\rm rep}/2,\alpha/(2C_{\rm rep})\}.
\]
All local proposals use the same primitive permutation list as the
tester, enlarging it by a fixed number of maps if necessary.

For the correction claim, let $d$ and $\delta_1$ be the outer A-code
degree and distance. Choose $\rho_{\rm raw}$ smaller than both
$\delta_{\rm raw}/3$ and $\delta_1/4$, and smaller than the fixed
uncertainty radius in Lemma~\ref{lem:outer-uncertainty}. An admissible
choice of $U^+$ is the set of centers having at least $(\delta_1/2)d$
neighbors in $U$ in a forward or inverse A list. Off this set all row
suggestions have fewer than half the local minimum distance in errors,
so the joint decoder recovers the reference exactly; the clock majority
and padding agree there too. If $u=|U|/|\Gamma|$, the spectral estimate
from that lemma gives
\[
 \frac{|U^+|}{|\Gamma|}
 \leq\frac{2\lambda^2u}{(\delta_1/2-u)^2}
 \leq\frac u2.
\]
Neither this support nor the proof refers to the values inside $U$.
For a requested-center cone with consistent repeated positions, extend
its initially supplied values arbitrarily to the whole word, setting
the reference outside $U$. The same support argument applies. This
proves the assertion even when separate cones choose different
extensions.

After $O(\log|\Gamma|)$ synchronous global rounds the uncertainty set
has cardinality less than one, giving exact recovery in polynomial
time. From a recovered logical singleton slot, read its coefficient
payload and evaluate the degree-$<h$ interpolant at $\omega^j$,
$0\leq j<h$, to obtain the corresponding logical bits. The repeated
clock gives $(c,\eta,j)$. The distance choice makes the promised
reference unique. On inputs outside the recovery promise the algorithm
has a total default; it may check the decoded label and its canonical
word by the global polynomial-time algorithms. Finally,
\eqref{oe:recipe-size-recurrence}--\eqref{oe:total-raw-size} and
\eqref{oe:fixed-interface-composition} give all asserted cost and
parameter-order statements.
\end{proof}
